\section{Mathematical Interpretations of the Direct Map}
\label{sec:geometry-notes}

The following three observations separate properties of a direct prediction rule from properties that would require evidence about a trained decoder.

\subsection{What can a decoder infer without a noise level?}

Fix the encoder and let $X$ be a clean training sample, $Z=E_\phi(X)$ its code, $A$ a random corruption strength, and $\widetilde{X}=C_A(X,\epsilon)$ the state presented to the decoder. Under squared-error reconstruction, an unrestricted optimal predictor is the conditional mean
\begin{equation}
  D^*(u,z)=\mathbb{E}[X\mid\widetilde{X}=u,Z=z].
  \label{eq:posterior-state}
\end{equation}
Let $\mu_\alpha(u,z)$ denote that mean when the decoder also observes $A=\alpha$. The law of total expectation gives
\begin{equation}
  D^*(u,z)=\int \mu_\alpha(u,z)\,
  P(d\alpha\mid u,z).
  \label{eq:latent-noise-level}
\end{equation}
The posterior law $P(d\alpha\mid u,z)$ weights corruption strengths by how well each explains the observed state and code. If several strengths are plausible, the optimal prediction combines their clean-state estimates in those proportions. This gives a precise interpretation of using one map across levels under squared error. The Pandora decoder uses a perceptual loss and the gencoder uses an $L_1$ loss; neither is trained to return the conditional mean in Eq.~\eqref{eq:posterior-state}.

\subsection{When do repeated calls repair a state?}

Write $D_z(u)=D_\theta(u,z)$ for the decoder with code $z$ held fixed. Suppose a desired reconstruction $u^*$ is a fixed point: $D_z(u^*)=u^*$. If a neighborhood $\mathcal{B}$ is invariant under $D_z$ and there is a constant $0\leq q<1$ such that
\begin{equation}
  \|D_z(u)-u^*\|\leq q\|u-u^*\|,
  \qquad u\in\mathcal{B},
  \label{eq:repair-contraction}
\end{equation}
then $\|u_k-u^*\|\leq q^k\|u_0-u^*\|$. A perturbation that leaves the state inside $\mathcal{B}$ can therefore be repaired by repeated calls to the decoder. The scalar $V(u)=\tfrac12\|u-u^*\|^2$ gives the basin illustration a precise reading: after one update, $V(D_z(u))\leq q^2 V(u)$. This is a sufficient local condition, not a consequence of the reconstruction objective. The bound on one multiplicative gate in the Pandora denoiser also gives no such bound for the whole network.

\subsection{In what sense is a direct update an Euler step?}

For any direct map $D_z$, define $\Delta_z(u)=D_z(u)-u$. Construct the autonomous ODE
\begin{equation}
  \frac{du}{d\tau}=\Delta_z(u).
\end{equation}
One Euler step of length one then gives $u_{k+1}=u_k+\Delta_z(u_k)=D_z(u_k)$. This algebraic identity holds for any state map. It says nothing about how accurately the discrete iterates follow a continuous trajectory, whether $\Delta_z$ is the gradient of a potential, or whether the trajectory follows a probability-flow ODE. In ReGAE the denoiser is trained on the target state itself.
